\section{The original four factors determine the order and the actual moments}
The receiving map in EQR is built from the original moments, so its
exceptional jet equations must retain their origin. We now evaluate
that origin using the four actual factors of FC13, with the original
period family NC1--5 and PCL1--16 \cite{NC,Cumulative,PCLOriginal} and the original Gamma
coefficient conventions FC20--24 \cite{WCF}. No period matrix, unit,
branch factor, marked root or original moment is replaced.

The finite four-factor test and the bound $v\le32$ were already
proved in the programme source PZ12 \cite{PeriodZeroOriginal}. The
calculation here retains that source's credit, evaluates its exact
conjugate-period and weighted-factor returns, and uses them in the
marked receiver below.

\subsection{Every factor coefficient and its full finite jet inverse}
Retain the ordered quartet $\rho_a$ of FC5, $Q_0(t)=t_1t_4-t_2t_3$,
the entire original $H$, and $D_5=\operatorname{diag}(\zeta,\zeta^2,
\zeta^3,\zeta^4)$. Define
\[
 B_0=\frac12\begin{pmatrix}0&0&0&1\\0&0&-1&0\\
 0&-1&0&0\\1&0&0&0\end{pmatrix},\quad
 M_j=H^{-1}D_5^{-j}H,\quad C_j=M_j^{\mathsf T}B_0M_j,
 \quad f_j(z)=t(z)^{\mathsf T}C_jt(z),\quad 1\le j\le4.
 \tag{OFM1}
\]
Thus $f_j=Q_j(Ht(z))$ exactly and $E_A=\prod_{j=1}^4f_j$.
Let $\alpha=(1,1,0,0)$, $\beta=(1,0,1,0)$ record the original
four labels. The full collected coefficients are
\[
 c_{j,uw}=\sum_{\substack{a,b\in\{1,2,3,4\}\\
       \alpha_a+\alpha_b=u,\ \beta_a+\beta_b=w}}(C_j)_{ab},
 \quad \lambda_{uw}=1+(2u-2)\delta+i(2w-2)\gamma,
 \quad f_j(z)=\sum_{u,w=0}^2c_{j,uw}e^{\lambda_{uw}z}.
 \tag{OFM2}
\]
Ordered pairs in this formula retain both off-diagonal matrix entries.
It follows by adding the two original exponents $\rho_a+\rho_b$,
so every sign, coefficient and original unit is present.

Order the nine pairs $(u,w)$ lexicographically, and set
$\mathcal V_{n,(u,w)}=\lambda_{uw}^{\,n}$ for $0\le n\le8$.
All nine exponents are distinct: equality of real parts gives $u=u'$
because $\delta>0$, and then equality of imaginary parts gives $w=w'$
because $\gamma>0$. Therefore
\[
 \det\mathcal V=\prod_{a<b}(\lambda_b-\lambda_a)\ne0,
 \quad a_{j,n}:=f_j^{(n)}(0)=\sum_{u,w}c_{j,uw}\lambda_{uw}^{\,n},
 \quad c_j=\mathcal V^{-1}(a_{j,0},\ldots,a_{j,8})^{\mathsf T}.
 \tag{OFM3}
\]
The determinant follows by subtracting columns in the Vandermonde
matrix. Its inverse is the adjugate divided by that exact determinant.
This supplies the complete finite-jet morphism and its inverse.

Each $f_j$ is nonzero on the admitted five-orbit domain. Indeed the
degree-two kernel of the exponential substitution is exactly
$\mathbb C Q_0$, as proved in FC12. If $f_j=0$, OFM1 would therefore
give $Q_j(Ht)=cQ_0(t)$ as polynomials. Substitution $t=H^{-1}X$
would make $Q_j(X)=cQ_0(H^{-1}X)$, contradicting the original
nonassociate-orbit condition. The invertible congruence in OFM1 also
excludes $c=0$. Consequently the finite test
\[
 n_j=\min\{0\le n\le8:a_{j,n}\ne0\},\quad
 v=n_1+n_2+n_3+n_4\le32,\quad
 \boxed{\mu_v=v!\prod_{j=1}^4\frac{a_{j,n_j}}{n_j!}\ne0}
 \tag{OFM4}
\]
always terminates and evaluates the actual original order and leading
moment. Multiplication of the four convergent Taylor series proves
both the order equality and the exact leading coefficient. It is the
original factorization, not a count of the collected degree-eight
support alone, that improves the earlier $81$-frequency bound $80$ to $32$.
For every $n$, including the four moments used in the actual receiver,
the complete formula is
\[
 \mu_n=\sum_{k_1+\cdots+k_4=n}
       \frac{n!}{k_1!k_2!k_3!k_4!}
       a_{1,k_1}a_{2,k_2}a_{3,k_3}a_{4,k_4},\qquad
 a_{j,k}=\sum_{u,w}c_{j,uw}\lambda_{uw}^{\,k}.
 \tag{OFM5}
\]
Thus no unspecified replacement moments are needed in EQR1.

\subsection{The exact conjugation relation of the original period family}
For this part use precisely NC1--4:
\[
 H(u)=D_u\mathcal R(u^{-1})V^{-1}U,\quad
 U=\operatorname{diag}(a,\bar a,\bar a,a),\quad
 V_{a,k}=w_a^k,\quad w_a=\rho_a-1/2,
 \quad a\ne0.
 \tag{OFM6}
\]
The scalar, diagonal Gamma-ray constants and branch powers are all
retained in the actual invertible $D_u$. Put
$R_u=\mathcal R(u^{-1})$ and $L_u=U^{-1}VR_u^{-1}$.
Since $D_u$ commutes with $D_5$, direct multiplication gives
\[
 H(u)^{-1}D_5^{-j}H(u)=L_uD_5^{-j}L_u^{-1}.
 \tag{OFM7}
\]
This equality proves the full diagonal cancellation in this particular
conjugated operator and restores it by OFM6; it does not replace the
original physical coordinates or their metrics.

Let $J$ exchange labels $1,2$ and $3,4$. The original quartet and
unit give $\bar V=JV$, $\bar U=JUJ$, and $Q_0(Jt)=-Q_0(t)$.
Every coefficient of NC4 is real for the original real $\delta,\gamma$.
Its convergent series therefore gives
$\overline{R_u}=R_{\bar u}$, including every coefficient, and hence
\[
 \overline{L_u}=J L_{\bar u},\qquad
 \overline{t(\bar z)}=Jt(z),\qquad
 \overline{L_u^{-1}t(\bar z)}=L_{\bar u}^{-1}t(z).
 \tag{OFM8}
\]
For a function $h$, write $h^\#(z)=\overline{h(\bar z)}$.
Using $\overline{D_5^{-j}}=D_5^j$ and OFM7--8 inside $Q_0$
proves the exact relation
\[
 f_{j,u}^{\#}(z)=-f_{5-j,\bar u}(z).
 \tag{OFM9}
\]
The sign is exactly the sign in $Q_0(Jt)$. In particular the
admitted orbit condition passes to the conjugate period, since
conjugation permutes its nonassociate quadrics and applies an
invertible coordinate map.

Define $g_u=f_{1,u}f_{2,u}$. Its full exponential coefficients are
$b_{uw}=\sum_{u_1+u_2=u,w_1+w_2=w}c_{1,u_1w_1}c_{2,u_2w_2}$,
with $0\le u,w\le4$, and its exponents are
$2+(2u-4)\delta+i(2w-4)\gamma$. The two minus signs in OFM9
cancel, giving
\[
 E_{A,u}(z)=g_u(z)g_{\bar u}^{\#}(z),\qquad
 \mu_n(u)=\sum_{k=0}^n\binom nk
          g_u^{(k)}(0)\overline{g_{\bar u}^{(n-k)}(0)}.
 \tag{OFM10}
\]
This is an exact identity also for nonreal admitted periods, and
retains the stated conjugate-period map. For real admitted $u$ it
becomes $E_A(x)=|g_u(x)|^2\ge0$ on the real axis. Put
$\nu=n_1+n_2\le16$ there. Then
\[
 n_4=n_1,\quad n_3=n_2,\qquad
 v=2\nu\le32,\qquad
 \boxed{\mu_v=\binom{2\nu}{\nu}|g_u^{(\nu)}(0)|^2>0},
 \qquad \mu_n\in\mathbb R\quad(n\ge0).
 \tag{OFM11}
\]
Orders are preserved by conjugate reflection; multiplying the two
Taylor series gives the displayed coefficient. Conjugating the
convolution sum interchanges $k$ and $n-k$, proving reality of every
moment. Neither higher-moment nonnegativity nor an adjoint identity
for the weighted conductor is inferred from this real-axis identity.
The positive-period instance of this positivity already follows from
PCL15--16: substitute $Z=e^{2\delta x}$ and $W=e^{2i\gamma x}$
in its exact slice to obtain
$E_A(x)=e^{(4-8\delta)x}W^{-4}A_{u^{-1}}(Z,W)$.
The additions here are the individual-factor identity OFM9 for the
conjugate period, the exact leading moment OFM11, and their use in
the original receiver and weighted four-factor map below.

\subsection{The full weighted factorization of the original conductor}
Let $D\ge0$ be the unchanged polynomial cutoff, and keep the original
$s$, Gamma mass $M=(2\pi)^{s/2}$, and weights
$\rho_m=\sqrt{m!/(s/2)_m}$. Set
$r_j=\sum_{k=j+1}^4n_k$ for $0\le j\le4$, so $r_0=v,r_4=0$.
For each original factor define its literal translation operator
\[
 T_jF(S)=\sum_{u,w=0}^2 c_{j,uw}F(S+\lambda_{uw}),\quad
 \mathcal H_j=\mathcal P_{D+r_j}/\mathcal P_{r_j-1},\quad
 c_j^{\rm centre}=a_{\rm amp}/2-j.
 \tag{OFM12}
\]
The quotient by $\mathcal P_{-1}$ means no quotient. Equip each
$\mathcal H_j$ with its original-type Gamma quotient norm at the
displayed centre, including the same mass $M$ and parameter $s$.
The four maps $T_j:\mathcal H_{j-1}\to\mathcal H_j$ are well-defined
and invertible: their Taylor formula is
$T_j S^N=\sum_{k\ge n_j}\binom Nk a_{j,k}S^{N-k}$.
It sends the old denominator to the next denominator and sends each
remaining leading degree $r_{j-1}+q$ to $r_j+q$ with nonzero
coefficient $\binom{r_{j-1}+q}{n_j}a_{j,n_j}$.
Distinct leading degrees prove both bijectivity and the full
triangular inverse by successive elimination.

The translations commute; multiplying their finite exponential
symbols or just expanding the four sums gives
\[
 T_A=T_4T_3T_2T_1:\mathcal H_0\longrightarrow\mathcal H_4.
 \tag{OFM13}
\]
Its domain, codomain and centres are exactly the original ones in
FC21, and no intermediate kernel is dropped.
Here are the full original-metric matrices. With
$\omega(t)=-i\arctan t$, set
\[
 G_j(t)=t^{-n_j}e^{-\omega(t)}f_j(\omega(t)),\quad
 (G_j)_0=(-i)^{n_j}a_{j,n_j}/n_j!,\quad
 A_j=\left[\operatorname{diag}(\rho_{r_j+q})_{q=0}^D\right]^{-1}
       T_D(G_j)\operatorname{diag}(\rho_{r_{j-1}+q})_{q=0}^D.
 \tag{OFM14}
\]
The coefficient formula follows from the same original orthogonal
polynomial generating series as FC22, with centre displacement one
and source/target degrees $r_{j-1}+q,r_j+q$. Explicitly the
generating-series translation multiplier is
$e^{-\omega(t)}\sum c_{j,uw}e^{\lambda_{uw}\omega(t)}$;
its first nonzero term has degree $n_j$. Extracting input and output
coefficients gives $(G_j)_{q-r}\rho_{r_{j-1}+q}/\rho_{r_j+r}$
for $r\le q$, and zero otherwise. This proves every matrix entry,
including the index offsets and original Gamma factors.

Products telescope all intermediate diagonal weights and satisfy
$\prod_jG_j=t^{-v}e^{-4\omega(t)}E_A(\omega(t))=g(t)$ from FC20.
Therefore the complete finite matrix identity and its inverse are
\[
 A_{D,s}=A_4A_3A_2A_1,\qquad
 A_{D,s}^{-1}=A_1^{-1}A_2^{-1}A_3^{-1}A_4^{-1}.
 \tag{OFM15}
\]
For every factor the inverse is fully explicit: if
$B_j(t)=\sum_{q=0}^D b_{j,q}t^q$, set
$b_{j,0}=(G_j)_0^{-1}$ and
$b_{j,q}=-(G_j)_0^{-1}\sum_{k=1}^q(G_j)_k b_{j,q-k}$.
Then $A_j^{-1}=\operatorname{diag}(\rho_{r_{j-1}+q})^{-1}
T_D(B_j)\operatorname{diag}(\rho_{r_j+q})$.
Truncated convolution proves both inverse products, so OFM15
reconstructs precisely the inverse in FC24 with no new assumptions.
For every exterior degree $0\le k\le D+1$, functoriality of the
exterior power and its operator norm gives
\[
 \|\wedge^k A_{D,s}^{-1}\|
   \le\prod_{j=1}^4\|\wedge^k A_j^{-1}\|,
 \quad
 \det A_{D,s}=
 \left(\prod_{j=1}^4\frac{(-i)^{n_j}a_{j,n_j}}{n_j!}\right)^{D+1}
 \prod_{q=0}^D\frac{\rho_{v+q}}{\rho_q}.
 \tag{OFM16}
\]
The equality is the unchanged original determinant because OFM4
identifies its parenthesis with $(-i)^v\mu_v/v!$. This is an
explicit factorization of the very inverse-exterior operator in
equation (11), with its original norms and every factor coefficient.

Finally take $D=3$ and the exact marking/interpolation maps
$\Phi_*,\Psi_*$ of FC26--32. Equation OFM13 gives all four original
ES--Fable columns, and every affine signed state, the same return:
\[
 T_4T_3T_2T_1\Phi_*q=\Psi_*q,\qquad
 M_*=C_\kappa B_{v,*}(V_\rho^{\mathsf T})^{-1}K^{-1},
 \quad (B_{v,*})_{r q}=\binom{v+q}{v+q-r}\mu_{v+q-r}
 \ (0\le r\le q\le3).
 \tag{OFM17}
\]
In this formula each moment is the evaluated finite sum OFM5, and
every $a_{j,k}$ is the finite original matrix expression OFM1--3.
On real admitted periods they additionally satisfy OFM11.
Thus the exceptional polynomial in EQR11 is connected to the actual
period and its exact factor jets by a complete finite coefficient
map. A fixed-old-quotient order jump still means the simultaneous
vanishing of the appropriate original moments as in FC25; the bound
and factorization here do not assert such a period exists or infer
RH from the separate nonlinear escape.
